\subsection{Evaluation Criteria}

The methodology requires evaluation criteria that are suitable for a research prototype and for the physical nature of the problem. Since the thesis does not claim full field validation of thermoelastic wave propagation in real geological media, the evaluation is framed as a comparative and consistency-oriented analysis rather than as final verification against complete real-world measurements.

The first criterion is physical consistency. The predicted response should be compatible with the main physical relationships introduced in Chapters 1 and 2. For example, changes in density, stiffness, thermal expansion, or thermal conductivity should produce qualitatively reasonable changes in the response. A model output that contradicts basic physical expectations without explanation should be treated with caution.

The second criterion is directional response. Since the thesis focuses on directional prediction, the response should be examined with respect to source--probe direction, propagation vector, distance, and azimuth. Even when the simplified mathematical formulation assumes isotropic effective properties, the framework must preserve the ability to describe and compare responses along different directions. This is especially important for interpreting layered, fractured, porous, or anisotropic geological media.

The third criterion is stability of outputs. A useful predictive component should not produce strongly erratic changes under small, physically minor changes in the input unless there is a clear physical reason for such behavior. Stability is important because thermoelastic interpretation depends on trends across materials, thermal conditions, and directions, rather than on isolated numerical outputs.

The fourth criterion is comparison between rock types. Sandstone, basalt, granite, limestone, and other representative geological materials may differ in density, elastic stiffness, thermal conductivity, heat capacity, and thermal expansion. The methodology should allow these differences to be reflected in the predicted response. The purpose is not to reduce geological behavior to a single material label, but to compare how different effective material parameters influence thermoelastic behavior.

The fifth criterion is comparison between models. Since several predictive components are considered, their outputs should be evaluated under comparable input conditions. This makes it possible to identify whether the models show similar trends, whether one model gives unstable or physically inconsistent responses, or whether some approaches are more interpretable for the selected prediction task.

The sixth criterion is interpretability. A prediction is more useful when it can be related back to physical variables such as temperature perturbation, displacement components, stiffness, density, or direction. Interpretability is particularly important in a thesis that is grounded in physical and geological formulation. The model output should support physical reasoning rather than remain an unexplained numerical value.

The final practical criterion is runtime and availability within the research prototype. Since the work is positioned as a prototype framework, a model component must be usable within the available system and should return responses in a practical manner. This criterion does not replace physical evaluation, but it is relevant for comparing model components in an implemented workflow.
